\section{改号、边镇与开元的转运：中心不能独自决定秩序}
\label{sec:wu-zhou-curated-chronicle}

670年大非川失利后，唐在青海与吐蕃之间的军政布局受到重压。\eventanchor{ST-670-DAFEICHUAN}{大非川之战（670）}战败不能化约为单一将领的失误：高原道路、补给距离、吐蕃的动员和唐朝的边军体制共同限定了可选行动。正史保存唐方的军事叙事，因而特别需要与吐蕃材料及高原考古的局部证据相互校读；这些材料也不能替我们复原每一支部队的兵数或意图。

690年武后改唐为周。\eventanchor{ST-690-ZHOU-FOUNDATION}{武周建国（690）}这一改号经过诏令、官僚请愿、宗室处置与佛教符瑞的语言，不是宫廷里的一次孤立宣告。695年契丹首领李尽忠、孙万荣起事，使河北道路、边军和朝廷的用人问题同时暴露。\eventanchor{ST-695-KHITAN-REBELLION}{契丹起事（695）}692年唐军恢复安西四镇的记载则表明，河西与西域并未因中原改号而停止成为资源与交涉的场域。\eventanchor{ST-692-ANXI-RECOVERY}{恢复安西四镇（692）}

\begin{debate}{复唐叙事能否抹平武周的行政连续？}
705年张柬之等迫武则天退位、恢复唐号。\eventanchor{ST-705-TANG-RESTORATION}{恢复唐号（705）}《旧唐书》《新唐书》与《资治通鉴》都在复唐后的正统语言中安排这一结果；年号墓志、官衔与制度条目提示人事和行政并非一夜归零。可确认的是政权名称与宫廷控制改变了，不能量出各地官员、寺院或家庭的认同。
\end{debate}

712年玄宗即位、713年平太平公主，重新收束了宫廷继承。\eventanchor{ST-712-XUANZONG-ACCESSION}{玄宗即位（712）}\eventanchor{ST-713-TAIPING}{平太平公主（713）}随后宇文融的括户财政与裴耀卿的转运方案，将国家能力落实到户籍、仓储和江淮水路上。\eventanchor{ST-721-YUWEN-RONG}{宇文融括户（721）}\eventanchor{ST-733-PEI-YAOQING}{裴耀卿转运（733）}它们不是“开元盛世”的装饰：多报一户、调一船粮，都会改变地方官、农户与漕运节点承担的成本。

\begin{sourcenote}{财政数字与制度文字的限度}
《唐会要》《通典》与两《唐书》可说明括户、转运的制度命名和朝廷争论；其数字与成效常由后来的编纂者概括。敦煌、吐鲁番文书能校验部分行政实践，仍不能据一地样本推出全国执行率。
\end{sourcenote}

开元的中枢整合与边地资源由此相连，也相互掣肘。751年怛罗斯战事使唐在中亚的推进受阻。\eventanchor{ST-751-TALAS}{怛罗斯战事（751）}读者可再沿\eventref{ST-751-AN-LUSHAN-COMMANDS}{安禄山兼领三镇（751）}继续追问：当边防、转运与将领任用集中时，这种整合会如何转变为危机。
